\chapter{Programming Language Implementation and Verification}
\label{ch:ch05}

\begin{objectives}
\begin{itemize}[leftmargin=*]
  \item Describe how the model is organized as a program and why it is written in JavaScript.
  \item Follow one breath through the code, from breath setup to the slow-state updates.
  \item Use the blind-simulation export to obtain reproducible numerical results.
  \item Distinguish verification (solving the equations correctly) from validation, and list the verification tests applied to the model.
\end{itemize}
\end{objectives}

\section{Choice of language}
\label{sec:impl-language}
VentiVerse is written in JavaScript and runs entirely inside a web browser, as a single HTML file that contains the model, the user interface and the graphics. JavaScript was chosen for reach rather than speed: a learner needs only a browser, the same file runs on a desktop, a tablet or a phone, and nothing has to be installed or licensed. Modern browser engines compile JavaScript to machine code as it runs, so the \SI{2}{\milli\second} loop of \cref{ch:ch03} executes in double-precision floating point much faster than real time. The same equations could be written in MATLAB, Python or C; the cost per simulated second would change, but the results would not, because the method and step are fixed (\cref{ch:ch04}).

\section{Program structure}
\label{sec:impl-structure}
The program separates three layers, shown in \cref{fig:impl-layers}. The \emph{model} layer holds the equations of Part III and Part IV and a single state object that carries every state variable from one step and one breath to the next. The \emph{scheduler} decides when breaths are computed: in interactive use, breaths are computed just ahead of the moving display, and in a blind run they are computed back to back as fast as possible. The \emph{interface} layer reads settings from the controls and draws waveforms, loops and trends; it never changes the model's equations. Keeping the layers apart is what allows the same model to be run by a learner moving a slider and by a script running a 12-hour validation night.

\begin{figure}[ht]
\centering
\begin{tikzpicture}[
  lay/.style={draw=bkblue,fill=bkblue!6,rounded corners=2pt,minimum width=11cm,minimum height=1.1cm,align=center,font=\sffamily\small},
  sig/.style={-{Latex[length=2mm]},thick,bkgray}]
\node[lay,draw=bkteal,fill=bkteal!6] (ui) at (0,2.6) {\textbf{Interface}: controls, waveform and loop displays, trends, reports};
\node[lay,draw=bkorange,fill=bkorange!8] (sch) at (0,1.3) {\textbf{Scheduler}: interactive (just ahead of display) or blind (back to back, CSV export)};
\node[lay] (mod) at (0,0) {\textbf{Model}: breath setup, \SI{2}{\milli\second} Euler loop, gas exchange, exact slow-state updates; one state object};
\draw[sig] ([xshift=-2cm]ui.south) -- node[left,font=\scriptsize]{settings} ([xshift=-2cm]sch.north);
\draw[sig] ([xshift=2cm]sch.north) -- node[right,font=\scriptsize]{breaths} ([xshift=2cm]ui.south);
\draw[sig] ([xshift=-2cm]sch.south) -- node[left,font=\scriptsize]{request next breath} ([xshift=-2cm]mod.north);
\draw[sig] ([xshift=2cm]mod.north) -- node[right,font=\scriptsize]{breath record} ([xshift=2cm]sch.south);
\end{tikzpicture}
\caption{The three layers of the program. Only the model layer changes the state of the simulated patient.}
\label{fig:impl-layers}
\end{figure}

\section{One breath through the code}
\label{sec:impl-breath}
Each breath is computed by one call to the breath builder, which follows the cycle of \cref{fig:num-loop}:
\begin{enumerate}[leftmargin=*]
  \item \textbf{Read settings.} The ventilator mode and its settings, the patient's mechanics and the inspired gas are read once, so that a setting changed during a breath takes effect at the next breath, as on a real ventilator.
  \item \textbf{Set the breath.} The controller computes the neural drive from the sensed carbon dioxide, oxygen saturation and pH, then the breath's duration, inspiratory time and muscle-pressure waveform (\cref{ch:ch07,ch:ch08}).
  \item \textbf{Integrate.} The Euler loop advances the compartment volumes and the airway pressure every \SI{2}{\milli\second}. Within the loop the ventilator logic checks for triggering and cycling, and the nonlinear terms (flow limitation, upper-airway collapse, tube resistance, leak) are evaluated at the current state. Pressure, flow, volume and muscle pressure are stored for display.
  \item \textbf{Summarize.} The breath's tidal volume, alveolar ventilation and dead space are computed and passed to gas exchange (\cref{ch:ch09,ch:ch10}).
  \item \textbf{Update the slow states.} PaCO\textsubscript{2}, the tissue carbon dioxide store, the chemoreceptor signals, bicarbonate and muscle fatigue are advanced over the breath duration by the exact updates of \cref{sec:num-exact}.
\end{enumerate}
The breath record produced by these steps contains everything the interface needs to draw the breath and everything a blind run writes to its output file.

\section{Blind simulation and data export}
\label{sec:impl-blind}
All numerical results in this book were obtained with the simulator's \emph{blind run}, which computes a chosen number of hours without drawing the display and exports the result as a comma-separated file. The file has three parts: a header that records the software version, the patient, the therapy, the mechanics, the sleep schedule and a complete snapshot of every setting; a waveform table at \SI{100}{\hertz} with airway pressure, flow, volume, muscle pressure, a simulated esophageal pressure and the breath phase; and a breath table with one row per breath containing more than 40 quantities, among them rate, tidal volume, minute and alveolar ventilation, peak and end-expiratory pressures, blood gases, pH, end-tidal carbon dioxide, dead-space fraction, cardiac output, muscle effort and fatigue, intrinsic PEEP and the occlusion pressure P\textsubscript{0.1}. Because the settings snapshot is stored in the file, any result can be reproduced by loading the same settings, and any figure in this book can be traced to the file it came from.

Scripted validation runs drive the same page through a headless browser: a script selects the phenotype, sets the patient's demographics, muscle strength and therapy, starts a blind run and saves the exported file. The model code run by the script is the same file a learner opens.

\begin{simnote}
The blind run is opened from the simulator's run menu. Set the duration, choose whether waveforms are exported, and optionally set a sleep period; the file downloads when the run ends.
\end{simnote}

\section{Breath-to-breath variability and repeat runs}
\label{sec:impl-random}
The model deliberately varies breathing from breath to breath (\cref{ch:ch08}), and sleep-stage transitions and arousals are stochastic (\cref{ch:ch11}). These use the browser's random-number generator without a fixed seed, so two runs with identical settings differ breath by breath, just as two nights of the same patient differ. Averages over many breaths are reproducible; individual breaths are not. Results that depend on rare events, such as the apnea--hypopnea index, are therefore reported in this book as the mean of repeat runs where the spread between runs matters.

\section{Verification}
\label{sec:impl-verification}
\emph{Verification} asks whether the program solves the model's equations correctly; \emph{validation} asks whether the equations represent real patients~\citep{oberkampf2010,asme_vv40}. This chapter deals only with verification. Validation is the subject of \cref{ch:ch14} and the disease chapters. The verification tests applied to the model are as follows.

\paragraph{Analytic solutions.} With a passive patient and a constant applied pressure, \cref{eq:num-eom} has an exponential solution with time constant $RC$ for each compartment. The Euler loop reproduces it within the error of \cref{tab:num-err}, and the pressure-control comparison of \cref{tab:solv-bench} confirms the full equation of motion with the pressure lag against a high-order reference solution.

\paragraph{Step size.} The full model was run for one hour of the chronic bronchitis patient used in the disease chapters, unsupported and on bilevel S/T 16/4, at the production step of \SI{2}{\milli\second} and at \SI{1}{\milli\second}. Because breathing varies at random, each \SI{2}{\milli\second} case was run twice, so that the difference between the two step sizes can be compared with the difference between two runs at the same step. \Cref{tab:impl-dt} gives the means over the second half-hour. Halving the step changed mean PaCO\textsubscript{2} by less than \SI{0.1}{\mmHg}, no more than the difference between two repeat runs at \SI{2}{\milli\second}, and changed tidal volume by at most \SI{1}{\milli\liter}. None of the outputs changed by a clinically meaningful amount.

\begin{table}[ht]
\centering
\caption{Step-size verification: one hour of the chronic bronchitis patient (\cref{ch:ch12}), means of breaths between 30 and 60 minutes. VT is the median tidal volume; RR is breaths per minute; the last column is carbon dioxide output, $\dot V_A\,\mathrm{PaCO_2}/0.863$, minus metabolic production, as a percentage of production.}
\label{tab:impl-dt}
\small
\setlength{\tabcolsep}{4pt}
\begin{tabular}{@{}llrrrrrrrr@{}}
\toprule
Therapy & Step & PaCO\textsubscript{2} & PaO\textsubscript{2} & SaO\textsubscript{2} & VT & RR & $\dot V_E$ & PEEPi & CO\textsubscript{2} out \\
 & & mmHg & mmHg & \% & mL & /min & L/min & cmH\textsubscript{2}O & \% \\
\midrule
Unsupported & 2 ms, run A & 57.3 & 53.6 & 83.7 & 335 & 19.1 & 6.36 & 3.03 & $-$3.0 \\
 & 2 ms, run B & 57.3 & 53.7 & 83.7 & 337 & 19.1 & 6.39 & 3.05 & $-$2.2 \\
 & 1 ms & 57.3 & 53.6 & 83.7 & 334 & 19.1 & 6.33 & 3.01 & $-$4.1 \\
\addlinespace
S/T 16/4 & 2 ms, run A & 49.7 & 58.9 & 89.4 & 492 & 12.0 & 5.91 & 2.60 & +6.5 \\
 & 2 ms, run B & 49.6 & 58.9 & 89.4 & 492 & 12.0 & 5.91 & 2.60 & +6.4 \\
 & 1 ms & 49.7 & 58.8 & 89.4 & 492 & 12.0 & 5.90 & 2.60 & +6.5 \\
\bottomrule
\end{tabular}
\end{table}


\paragraph{Carbon dioxide conservation.} At steady state, the carbon dioxide eliminated by alveolar ventilation, $\dot V_A \cdot \mathrm{PaCO_2}/0.863$, must equal the carbon dioxide produced. The last column of \cref{tab:impl-dt} gives the difference over the second half-hour as a percentage of metabolic production ($\dot V_{\mathrm{O_2}} \times$ respiratory exchange ratio). Unsupported, output was 2\% to 4\% below production in the three runs, while on S/T it was about 6\% above production. Neither run had fully reached steady state within the hour. On S/T, which started at time zero, PaCO\textsubscript{2} fell by about \SI{8}{\mmHg} and the body stores were still unloading carbon dioxide; the smaller deficit in the unsupported runs is consistent with the stores still filling from the starting blood gas. The balance varied by up to 1.8 percentage points between runs, about as much between repeat runs at \SI{2}{\milli\second} as between step sizes. \Cref{ch:ch10} examines the balance over longer runs.

\paragraph{Regression runs.} Before any change to the model is released, the scenarios of the disease chapters are rerun and compared with their previous results, so that a change made for one disease is seen in every other disease it affects.

\begin{summarybox}
The model runs as a single JavaScript program in the browser, with the model separated from the scheduler and the interface. Every result in this book comes from a blind run whose output file stores the settings that produced it. Verification tests show that the numerical error of the \SI{2}{\milli\second} Euler step is far smaller than the model's own breath-to-breath variability, and that carbon dioxide is conserved.
\end{summarybox}
